% Section 15.7, translated from sv/chapters/15/summary.tex.

\section{Summary}
\label{c15:sec:sammanfattning}

\begin{booktable}{@{}p{12.5em}L@{}}
  \thead{Concept} & \thead{Formula} \\
  \midrule
  Rectangular form & $z = x + jy$ \\
  Magnitude & $\abs{z} = \sqrt{x^2 + y^2}$ \\
  Phase angle & $\delta = \arctan(y/x)$ + correction \\
  Polar form & $z = \abs{z}\angle\delta$ \\
  Polar to rectangular form & $x = \abs{z}\cos\delta$, \ $y = \abs{z}\sin\delta$ \\
  Addition & $(x_1 + x_2) + j(y_1 + y_2)$ \\
\end{booktable}

\needspace{6\baselineskip}
\subsection*{What you should be able to do}
After this chapter, you should be able to:
\begin{itemize}
  \item Understand what complex numbers are and how they can be represented in the complex plane.
  \item Carry out simple calculations with complex numbers in rectangular form.
  \item Convert between rectangular and polar form, and give the magnitude and phase angle.
\end{itemize}

\testadigsjalv
\begin{enumerate}
  \item Convert $z = 3 + j4$ to polar form. Give the magnitude and the phase angle in radians.
  \item Calculate $(2 + j3)(1 - j)$ in rectangular form.
\end{enumerate}

\needspace{6\baselineskip}
\subsection*{Next chapter}
\Kapref{16} continues with complex numbers: Euler's formula, and complex numbers for vector
calculations.
